\documentclass[11pt,a4paper]{article}

\usepackage[T1]{fontenc}
\usepackage{helvet}
\renewcommand{\familydefault}{\sfdefault}
\renewcommand{\ttdefault}{\sfdefault}
\usepackage[top=2cm, bottom=2cm, left=2cm, right=2cm]{geometry}
\usepackage{booktabs}
\usepackage{xcolor}
\usepackage{parskip}
\usepackage{float}

\newcommand{\timecode}[1]{\textbf{\textcolor{blue}{[#1]}}}
\newcommand{\cue}[1]{\textit{\textcolor{gray}{$\triangleright$ #1}}}

\title{\textbf{COMP0240 CW1 --- Video Narration Script}\\
\large 4-Minute Submission Video}
\author{Ali Shihab}
\date{April 2026}

\begin{document}
\maketitle

\noindent\textbf{Runtime:} 4:00 limit. \textbf{Narration:} 392 words at \textbf{100\,wpm} $=$ 3:55, leaving 5\,s.\\
\textbf{Rubric:} Clarity 20\% \quad \textbf{Successful runs 60\%} \quad Performance analysis 20\%

\medskip
\noindent At 100\,wpm there is roughly a second of silence for every two spoken words, so the footage carries long stretches alone. Let it. The gaps between the lines below are where the drone should be visibly flying, not where you should speak faster.

\section*{Terminal layout and commands}

Three terminals, all inside the XRDP desktop (the camera window needs a real
display). \textbf{T1} runs the simulator, \textbf{T2} runs the mission, and
\textbf{T3} runs the camera viewer. Each begins with:

\begin{quote}\ttfamily
cd \textasciitilde/mission\_planning\_ws/src/challenge\_mission\_planning\\
source setup.bash
\end{quote}

\noindent\textbf{T1 --- simulator.} Relaunch between scenarios, and run
\texttt{./stop.bash} first so the previous world is fully torn down:
\begin{quote}\ttfamily
./stop.bash \&\& sleep 5\\
./launch\_as2.bash -s scenarios/scenario\emph{N}.yaml
\end{quote}

\noindent\textbf{T3 --- camera viewer} (start once the drone has spawned, leave running):
\begin{quote}\ttfamily
python3 aruco\_viewer.py --scale 0.6 --save-dir aruco\_captures
\end{quote}

\noindent\textbf{T2 --- the mission.} Every run uses the same parameters as the
72-run dataset, so what is on screen matches the report. Only
\texttt{--ordering} and \texttt{--local\_planner} change. Define this once per
terminal:
\begin{quote}\ttfamily
\footnotesize
MP="--namespace drone0 --use\_sim\_time --speed 3.0 --sample\_hz 50 \textbackslash\\
\hspace*{2em}--obstacle\_inflation\_m 0.4 --grid\_resolution 0.5 --goto\_timeout\_s 60.0 \textbackslash\\
\hspace*{2em}--stuck\_timeout\_s 20.0 --dwell\_s 0.2 --final\_pos\_tolerance\_m 0.75 \textbackslash\\
\hspace*{2em}--yaw\_tolerance\_deg 30.0 --return\_to\_start \textbackslash\\
\hspace*{2em}--aruco\_topic /drone0/sensor\_measurements/hd\_camera/image\_raw"
\end{quote}

\begin{table}[H]
\centering\small
\begin{tabular}{llll}
\toprule
\textbf{Segment} & \textbf{Sim (T1)} & \textbf{Mission (T2)} & \textbf{Shows} \\
\midrule
Sc1, 0:33 & \texttt{scenario1.yaml} & \texttt{\$MP --ordering input --local\_planner astar}   & file order \\
          &                         & \texttt{\$MP --ordering nn\_2opt --local\_planner astar} & 2-opt \\
\midrule
Sc2, 1:00 & \texttt{scenario2.yaml} & \texttt{\$MP --ordering input --local\_planner straight} & \textbf{fails, 0\%} \\
          &                         & \texttt{\$MP --ordering input --local\_planner astar}    & completes \\
\midrule
Sc3, 1:35 & \texttt{scenario3.yaml} & \texttt{\$MP --ordering input --local\_planner astar}    & 226\,m \\
          &                         & \texttt{\$MP --ordering nn\_2opt --local\_planner astar}  & 119\,m \\
\midrule
Sc4, 2:02 & \texttt{scenario4.yaml} & \texttt{\$MP --ordering nn\_2opt --local\_planner straight} & \textbf{fails, 0\%} \\
          &                         & \texttt{\$MP --ordering nn\_2opt --local\_planner astar}    & completes \\
\bottomrule
\end{tabular}
\end{table}

\noindent Prefix each with \texttt{python3 mission\_scenario.py scenarios/scenarioN.yaml}.
Each pair varies exactly one factor, matching the argument in the narration.
Scenario~2 and~4 straight-line runs are expected to fail --- that is the footage
the argument depends on, so record them rather than retrying until they pass.

\medskip
\noindent\textbf{Recording requirements}
\begin{itemize}
  \item The 60\%-weighted criterion asks for \emph{``capturing ArUco marker images, providing clear visual proof''}. A terminal banner is not an image. Show the \textbf{camera feed} with the marker outlined and its ID drawn on (\texttt{cv2.aruco.drawDetectedMarkers}).
  \item \textbf{Layout:} Gazebo left ($\sim$60\%), the \texttt{aruco\_viewer} window top-right, and terminal \textbf{T2} bottom-right, 1920$\times$1080. T2 is the mission log: \texttt{mission\_scenario.py} prints arm, takeoff, every viewpoint result and each \texttt{Detected IDs} line to its own terminal (also tee'd to \texttt{runs/<run>/stdout.log}). The tmux session in T1 shows only simulator and ROS node output, which is not worth screen space.
  \item Scenario~3 at 2$\times$ with an on-screen caption saying so.
  \item The Scenario~2 and~4 baseline failures carry the argument. Show the drone driving into the obstacle and the watchdog giving up on the leg. The report calls strikes \emph{inferred} from altitude traces, not measured, so describe what is on screen and claim no collision sensor.
\end{itemize}

\bigskip\hrule\bigskip

\section*{Segment 1 --- Problem and Approach (0:00--0:33)}

\cue{Gazebo wide shot of Scenario~2: drone on pad, five obstacle blocks, ArUco markers on stands}

\timecode{0:00}
``The task is to fly a drone around a set of inspection viewpoints, photograph an ArUco marker at each one, and avoid every obstacle.

\timecode{0:12}
There are two decisions. What order to visit the viewpoints in, and how to route between them.

\timecode{0:20}
For ordering I use nearest-neighbour with 2-opt. For routing, either A-star on a voxel grid, or RRT-star sampling in continuous space.

\timecode{0:32}
Rather than pick a configuration and test it, I crossed both factors: two orderings, three planners, four scenarios, three repeats. Seventy-two runs.''

\bigskip\hrule\bigskip

\section*{Segment 2 --- Demonstration (0:33--2:30)}

\subsection*{Scenario 1 --- 10 viewpoints, 1 obstacle (0:33--1:00)}

\cue{Split-screen, both running A*: file ordering left, NN+2-opt right. Camera inset showing marker detections. Let it run.}

\timecode{0:33}
``Scenario 1. Same planner on both sides. Only the ordering differs.

\timecode{0:42}
On the left, file order. The path keeps crossing itself. On the right, 2-opt. No crossings.

\timecode{0:52}
That one change takes the path from 159 metres to 96.''

\subsection*{Scenario 2 --- 10 viewpoints, 5 obstacles (1:00--1:35)}

\cue{Baseline clip first: drone flying into an obstacle. Caption ``BASELINE --- OBSTACLE STRIKE''. Hold it. Then cut to A*.}

\timecode{1:00}
``Scenario 2 adds five obstacles. This is where flying straight breaks down.

\timecode{1:10}
It heads directly for a viewpoint sitting behind an obstacle, collides, and the watchdog abandons that leg. Over three runs it reached about half the viewpoints. It never once finished.

\timecode{1:26}
A-star routes around all five and completes every run. All ten markers confirmed.''

\subsection*{Scenario 3 --- 20 viewpoints, 1 obstacle (1:35--2:02)}

\cue{Split-screen at 2$\times$, ``2$\times$ SPEED'' caption}

\timecode{1:35}
``Scenario 3 doubles the viewpoints to twenty. This is the strongest case for the tour optimiser.

\timecode{1:47}
File order flies 226 metres. Two-opt flies 119. The tour drops from 242 seconds to 163.''

\subsection*{Scenario 4 --- 5 viewpoints, 10 obstacles (2:02--2:30)}

\cue{Baseline failing, then A* threading the obstacle field}

\timecode{2:02}
``Scenario 4 inverts it. Five viewpoints, ten obstacles.

\timecode{2:10}
With so few viewpoints there is nothing to reorder, so ordering buys nothing here. Avoidance is everything.

\timecode{2:20}
The baseline reaches two thirds of the viewpoints and never finishes. A-star completes all three runs. Its 73 seconds beats the baseline's 77, but that baseline never finished, so the comparison does not hold.''

\bigskip\hrule\bigskip

\section*{Segment 3 --- Results (2:30--3:30)}

\cue{Full-screen \texttt{report\_figs/fig\_reliability.png}, then the table below}

\timecode{2:30}
``Two findings, pulling in opposite directions.

\timecode{2:36}
Ordering drives efficiency. With the planner held fixed, 2-opt cuts path length by 38 to 50 per cent, and cuts tour time in every comparison where both configurations finished.

\timecode{2:52}
The obstacle-aware planners push efficiency the other way. A-star adds up to 40 seconds.

\timecode{3:02}
Now look at success. Straight-line flight hits zero per cent in two conditions. A-star completes all twenty-four of its runs. Those extra seconds are buying reliability.

\timecode{3:19}
One caution. Below a hundred per cent success, the times and distances mean nothing, because the drone simply did less work.''

\bigskip
\begin{table}[H]
\centering\small
\begin{tabular}{llccc}
\toprule
\textbf{Effect} & \textbf{Comparison (one factor varied)} & \textbf{Tour (s)} & \textbf{Path (m)} & \textbf{Success} \\
\midrule
\textbf{Ordering} & Sc3 \; File order / A* & 241.5 & 226.2 & 100\% \\
\emph{planner fixed} & Sc3 \; NN+2-opt / A* & \textbf{162.7} & \textbf{118.8} & 100\% \\
\midrule
\textbf{Planner} & Sc4 \; NN+2-opt / Straight & 77.3 & --- & \textbf{0\%} \\
\emph{ordering fixed} & Sc4 \; NN+2-opt / A* & 73.4 & 70.7 & \textbf{100\%} \\
\bottomrule
\multicolumn{5}{l}{\footnotesize Each pair varies one factor. Dashes: not comparable below 100\% success.}
\end{tabular}
\end{table}

\bigskip\hrule\bigskip

\section*{Segment 4 --- Conclusions (3:30--3:55)}

\cue{Drone landing at origin, or the four-scenario path figure}

\timecode{3:30}
``So the two components do different jobs. Ordering makes the mission efficient. Obstacle-aware planning gets it finished.

\timecode{3:42}
Tune on time alone and you would pick the configuration that leaves the survey incomplete. Thanks for watching.''

\bigskip\hrule\bigskip

\section*{Timing budget at 100\,wpm}

\begin{table}[H]
\centering\small
\begin{tabular}{lccl}
\toprule
\textbf{Segment} & \textbf{Words} & \textbf{Time} & \textbf{Rubric served} \\
\midrule
1 Problem and approach & 55  & 0:33 & Clarity (20\%) \\
2 Demonstration        & 197 & 1:57 & \textbf{Successful runs (60\%)} \\
3 Results              & 105 & 1:03 & Analysis (20\%) \\
4 Conclusions          & 35  & 0:21 & Clarity (20\%) \\
\midrule
\textbf{Total}         & \textbf{392} & \textbf{3:54} & 6\,s spare \\
\bottomrule
\end{tabular}
\end{table}

\noindent Segment 2 takes 50\% of the runtime for a 60\%-weighted criterion, and at this pace most of it is footage rather than narration --- which is what that criterion rewards. If you run long, cut the caution at 3:19 ($\approx$14\,s). Never cut Segment 2.

\end{document}
